\section{拓扑的刚性：耦合密度、相变与结构记忆}

\subsection{网络耦合密度与虚拟质量}

有些事实只能由过程给出，这一点在时间维度上已经清楚：候选要进入材料、身体和组织，隐藏的后果才会逐步显现。但过程有一个前提：候选必须能被送入过程，目标必须能被系统接住。资源势垒讨论的是改变现实需要付出什么代价，时间刚性讨论的是现实过程不能被任意压缩，拓扑耦合讨论的则是另一个更深的问题：即使一个目标状态已经被找到，即使资源和时间也可以投入，系统是否存在一条能够抵达这个目标状态的路径？这就是拓扑耦合要处理的问题。

\subhead{从目标到路径}

要看清这个问题，先要改变对系统的描述方式。现实系统不是一组孤立变量的集合，而是由变量之间的关系构成的网络。道路与道路相连，设备与设备相连，部门与部门相连，规则与规则相连，人的行为与制度预期相连，技术接口与历史系统相连。一个元素的意义，不只由它自身决定，也由它所在的位置决定；改变一个元素，也不只是改变它自身，而是在改变它与周围结构的关系。因此，在复杂系统中，“局部变好”和“系统变好”之间没有必然通道。一个局部节点可以被优化，一个局部流程可以被加速，一个局部模型可以更准确，但这些改进能否传导为系统状态改变，取决于它们在网络中的位置、连接方式和边界条件。局部优化是点上的改进，系统改变是路径上的转移。二者之间隔着拓扑结构。

连接密度的作用可以先从记账方式上看。一个节点自身的规模决定了改动它本身需要多少工作量，这是它自己的部分；而改动它时必须一起调整的关系集合，构成了附加在它上面的另一部分。耦合密度上升，附加部分就变大，节点在行为上越来越不像一个可以单独推动的物体，而越来越像一个带着额外负载的物体：同样的推动力，得到的位移越来越小。把这部分附加的改动成本单独记账，可以把它称为节点在这个网络中的虚拟质量。这个量不是物理质量，它没有固定的单位，只能在同一系统的不同位置之间比较，它的用处是回答一个具体问题：为什么同一个改动，在这个位置是一天的工作量，在另一个位置是半年的项目。

这个量在网络科学里有对应的结构。随机图与规则格点给出的是度分布很窄的网络，节点之间的连接数差别不大；而真实的技术网络、组织网络和交通网络往往是度分布很宽的，少数节点承载了大量连接。Stanley Milgram 在 1967 年做的传递信件实验发现，任意两个美国人之间的中介链长度大约在六步以内\cite{milgram1967smallworld,travers1969experimental}，Duncan Watts 和 Steven Strogatz 在 1998 年给出了这种短路径与高聚集系数可以共存的网络模型\cite{watts1998collective}。这些结果说明两件事：其一，路径很短，任何局部扰动都能在很少的跳数内到达远处；其二，连接不是均匀分布的，改动代价集中在少数度很高的位置上。虚拟质量大的节点，正是那些度很高的节点，它们自身可能只是一个字段、一个接口、一张表，但它们连着的是整片系统。

技术接口与组织接口之间还存在一种对应关系，这一点在软件工程里被反复观察到：设计系统的组织结构，最终会以接口的形式复刻到系统结构里，不同团队之间的沟通边界往往会成为模块边界\cite{conway1968how}。这条经验在现实系统中表现为一种双重耦合，技术上的依赖关系和组织上的协调关系互相加强：一个接口之所以难以改动，不只是因为下游系统多，也因为改动它需要跨过若干部门的职责边界。判断一个节点的虚拟质量，因此不能只看它的连接条数，还要看这些连接是否跨越了组织边界、是否涉及外部承诺、是否已经写进了合同和考核。

\subhead{目标状态与状态转移}

路径是否存在，先要看被提出的那个状态是在什么空间里成立的。AI很擅长生成目标状态：它可以提出一个更高效的城市交通方案，一个更优的医院分诊流程，一个更低库存的供应链结构，一个更精准的教学评价模型，一个更自动化的工厂调度系统。可是，模型中生成的目标状态，只说明这个状态在表示空间中可以被描述，并不说明现实系统能够以低成本、低风险、稳定的方式抵达它。目标状态可描述，不等于状态转移可完成。

很多技术乐观主义的错误，正来自这里。它把“知道更优状态”误认为“系统能够进入更优状态”。但现实系统从来不是任意可改的。城市已经有道路、地铁、楼宇、管网、产权、居民习惯和财政约束；医院已经有科室结构、病历系统、医保支付、医生责任和患者信任；企业已经有旧软件、旧流程、部门利益、绩效体系和数据格式；教育系统已经有考试、升学、教师评价、家庭预期和社会筛选功能。一个方案进入这些系统时，不是在空白平面上画线，而是在已经存在的拓扑结构中寻找通道。所以，拓扑耦合的第一性含义是：现实系统的状态空间不是均匀开放的。某些状态虽然可以被想象，却没有可承受路径；某些方案虽然局部有效，却不能穿过系统接口；某些改进虽然在模型中成立，却会在真实结构中被吸收、延迟、扭曲或反向放大。技术落地的难点，不只是有没有一个好解，而是现实系统是否允许这个解成为可达状态。

这可以解释许多常见现象。一个AI诊断模型在测试集上准确率很高，但进入医院之后，它面对的不是抽象图像，而是医生工作流、患者病史、设备差异、医保规则、误诊责任和临床信任，模型的局部性能并不能直接改写医疗系统。一个交通优化算法在仿真中降低拥堵，但真实城市中，司机会调整路径，商圈会改变流量，施工会打断计划，执法方式会影响行为，公共交通和私人出行会相互反馈，一个看似更优的交通状态在系统中不一定有稳定路径。

\subhead{被系统吸收的改进}

路径不稳定，最常见的表现是改进在系统内部被吸收。一个企业引入AI助手，可能让每个员工生成文档更快，但组织整体未必更有效率。因为组织效率不只由文档生成速度决定，还由判断质量、沟通链条、审批结构、责任分配和目标一致性决定。如果每个人都更快地产生材料，而审核、协调和决策机制没有改变，系统可能只是从“生成瓶颈”转向“筛选瓶颈”。局部速度提高，整体状态未必改善。

\key{系统改变不是局部能力的简单求和。}一个复杂系统的整体状态，往往由关系结构决定，而不是由单个部件的性能决定\cite{boccaletti2006complex}。就像一张网，某个节点变得更强，并不意味着整张网能移动到新的形状；它必须通过线与其他节点传递力量，如果连接方式不允许，局部力量会被分散、抵消或锁住。更麻烦的是，这类阻力经常被误读为人的保守或对变化的反对。拓扑耦合给出的解释不同：现实系统中的许多阻力不是来自“反对变化”，而是来自“关系不能立即重排”。一项技术进入系统后，往往会改变谁负责、谁受益、谁承担风险、谁需要学习、谁的流程被打断、谁的权力被重分配。技术看似是工具，实际会触碰结构；结构越复杂，技术越不可能只是一个中性插件。

这类吸收有一个可以被观察到的形状：改进越靠下游，越容易被上游的固定节奏压平。一个车间把单件加工时间压缩了三分之一，如果装配线的节拍、物料的到货周期和检验的批次规则都没有变化，压缩出来的时间只会变成工序之间的等待，产出并不会等比上升。反过来，一项改进如果落在决定节拍的那一环上，即使幅度不大，也会整条链传导下去。判断一项改进能否成为系统改变，可以先找系统的节拍源在哪里，再看这项改进是否触及了它。

\subhead{结构位置与黏滞分布}

同一项技术在不同的位置上会有完全不同的命运。很多AI应用会经历“演示容易，部署困难；试点顺利，规模化缓慢”的路径。演示通常发生在低耦合环境中：任务明确，输入干净，接口预设，责任悬置，失败可以忽略；试点也常常选择最配合、最标准、最容易成功的局部场景。但规模化部署必须面对真实系统的异质性：不同机构的数据结构不同，人员能力不同，流程不同，利益结构不同，合规要求不同，责任边界不同，技术从样板进入系统时，拓扑成本才真正显现。

因此，判断一项AI技术的现实潜力，不能只问它能不能完成任务，还要问它嵌入什么结构。低耦合位置更容易自动化，高耦合位置更容易被限制。一个广告文案生成器、一个会议纪要工具、一个客服初筛系统，通常比医疗诊断、自动驾驶、教育评价、司法裁判、能源调度更容易落地，并不是因为前者的模型必然更强，而是因为前者所处的结构位置更轻，责任后果更小，失败更容易回滚。同一种能力，在不同结构位置上会表现出不同黏滞。同样是文本生成，用于营销草稿时，它主要连接传播效果；用于病历报告时，它连接诊疗责任；用于法律裁判时，它连接权利、义务和强制力。任务形式相似，拓扑位置不同，现实阻力也完全不同。

位置的差别还可以从接口数量上看出规模。一个只被内部少数脚本调用的函数，改动风险由作者自己承担；一个被上百个下游系统按固定格式调用的接口，改动风险由所有调用方共同承担，任何一次字段含义的调整都要走完通知、兼容、并行运行和旧版本退场的全过程。前者的改动以小时计，后者的改动以季度计。同样是几行代码，虚拟质量相差几个数量级，差别不在代码，而在它连着多少东西。

\subhead{改变连接方式的技术}

如果阻力来自位置，那么改变的入口就不在能力本身。真正改变系统的技术，往往不是单点能力最强的技术，而是能够改变连接方式的技术。智能手机改变世界，不只是因为硬件强，而是因为它重组了移动网络、应用商店、支付系统、开发者生态和用户行为；电动车扩散，不只是因为电机效率高，而是因为电池供应链、充电网络、政策、维修体系和消费认知逐渐成形。一个技术能否成为系统力量，取决于它能否找到或创造自己的生态位。所以，技术落地并不是把一个工具塞进现实，而是让现实系统为这个工具形成新的接口。接口清楚，责任明确，反馈闭合，失败可控，技术就容易进入；接口混乱，责任不清，反馈延迟，失败不可回滚，技术就会停留在演示和试点。拓扑耦合不是说系统不能被改变，它说的是系统改变需要结构通道：路径不是由目标状态自动给出的，而是由连接结构、历史路径、制度接口和行为模式共同打开或关闭的；一个技术能否进入现实，不只取决于答案是否清楚，还取决于答案是否能在系统中找到承接位置。

到这里可以回答开头那个问题：路径是否存在，取决于结构通道，而不取决于目标是否清楚。资源势垒告诉我们，现实改变需要支付物理账本；时间刚性告诉我们，现实改变需要经历过程；拓扑耦合告诉我们，现实改变需要穿过结构。没有资源，系统推不动；没有时间，过程展不开；没有路径，目标到不了。技术落地不是把答案放进现实，而是让答案找到结构通道。

但一条通道被打开，只是第一步。沿一条通道改变一个节点时，被牵动的往往不只是一处：接口、数据、权限、培训、责任和预期都要跟着移动，其中一部分必须同步移动，失败之后不一定能退回。这些成本决定了同样的改进为什么在不同的系统里快慢悬殊，也决定了一项技术进入系统之后，究竟会被吸收为局部效率，还是会把整个依赖网络一起拖入调整。

\subsection{拓扑相变：渗流阈值与凝胶化}

协同移动成本上升到某个区域之后，系统的响应方式会改变。复杂系统的困难，常常不是逐渐增加的，而是在某个区域突然改变性质。一个系统增加少量连接，可能只是提高协同；再增加一些连接，可能只是提高协调成本。但当连接密度、反馈回路和依赖范围达到某个临界区域之后，系统会表现出一种新的状态：它不再像松散部件那样可以局部调整，而更像一张贯通整体的网络。局部变化不再停留在局部，而会沿着连接传到远处；局部错误也不再容易隔离，而可能进入全系统的反馈回路。这就是“凝胶化系统”的基本含义。

这里的凝胶化不是严格的物理等同，而是一种结构类比。真正的凝胶化发生在聚合物、胶体、材料网络中：分散的单元在某个条件下突然形成贯通结构，系统从可以流动的状态变成带有整体支撑的状态。本章借用这一图像，不是要把社会系统、组织系统、交通系统和智能系统硬套成材料物理，而是要说明一种共同机制：\key{连接一旦越过某个区域，系统不再只是“复杂一点”，而会改变扰动传播的相态。}它从局部可吸收，变成全局可贯通；从局部可调整，变成整体性牵连；从柔软的试错空间，变成一张绷紧的网。

真实的凝胶化过程本身也有值得对照的地方，因为这个类比的来源并不是隐喻。高分子溶液在交联度上升时，黏度会先缓慢增加，随后在某个交联密度附近急剧上升，溶胶部分失去流动性，形成贯穿整个体系的网络；材料学把这个位置称为凝胶点，它对应的不是某个局部变硬，而是连通性第一次贯通到整个样品。凝胶点的位置可以用官能度与反应程度估算，而不是靠观察变得多稠来判断，这与后面要讨论的渗流阈值是同一类数学结构。真正需要记住的是这个位置的性质：在它之前，体系可以被局部搅动；在它之后，任何一次局部扰动都要由整个网络分担。

\subhead{渗流阈值与相互依赖}

这个区域有多宽，阈值在哪里，是可以用模型算出来的。渗流理论提供了一个朴素而有力的图像：一个网络在连接较少时，节点之间只是形成许多小的连通团簇，扰动很难跨越整个系统；当连接密度接近临界阈值时，小团簇开始合并；一旦越过阈值，网络中会出现巨连通分量，局部扰动就有可能穿过系统\cite{broadbent1957,erdos1959}。这一图像的价值不在于把现实世界简化为随机图，而在于提醒我们：连接增加带来的不是线性变化，而可能是相变。系统不是“越来越复杂一点”，而是从一种扰动传播机制跳到另一种扰动传播机制。

阈值本身可以用平均度写出。设 $\langle k\rangle$ 为网络的平均度，即每个节点平均拥有的连接数，则巨连通分量出现的条件是
\begin{equation}
\langle k\rangle_c = 1 .
\end{equation}
对Erd\H{o}s--R\'enyi随机图，平均度越过1时出现巨分量的渐近阈值。这里涉及随机分支的扩展，而非“每个节点都有一个连接就一定贯通”。$G(n,p)$的平均度为$(n-1)p$，对应的阈值密度约为$1/n$。实际依赖网络需要检查方向、相关性、权重与聚团。

在普通渗流的相应无限系统极限中，序参量可以在阈值以上连续地从零增长。有限网络表现为跨越概率和团簇规模的变化。设$p$为占据概率，$p_c$为阈值，$P_\infty$为巨分量中的比例，阈值附近可出现如下尺度形式：
\begin{equation}
P_\infty(p)\ \propto\ (p-p_c)^{\beta},\qquad p\to p_c^{+},
\end{equation}
这里 $\beta$ 是临界指数，$P_\infty$ 的取值在0与1之间，这个式子只在系统规模趋于无限时成立。$p_c$ 的取值依赖网络类型与维度：二维方格上的点渗流阈值约为 $p_c\approx 0.5927$，即每个格点被随机占据的比例要超过约59\%才会出现贯通分量；同一方格上的键渗流阈值是0.5；一维链上 $p_c=1$，必须全部占满才可能贯通。有限系统里没有精确的开关点，只有在阈值附近迅速变陡的过渡区，系统规模越小，过渡区越宽。

阈值的具体数值对网络类型很敏感，这一点在讨论现实系统时比 $0.5927$ 本身更重要。二维方格的坐标数是4，三维简单立方格子的点渗流阈值约为0.31，配位数越高、维度越大，贯通所需的占据比例越低；而在度分布很宽的网络上，如果存在少数度极高的节点，随机移除节点的阈值会趋近于零，也就是说，只要枢纽节点还在，随机损失大量普通节点也难以破坏整体连通\cite{albert2000error}。同一个网络对随机扰动的耐受和对蓄意攻击的脆弱并存，破坏它的方式不是均匀地增加压力，而是准确移除少数关键位置。渗流模型在这里给出的判断是分层的：阈值的高低取决于结构的形状，而不只是一个连接总数。

\begin{figure}[ht]
\centering
% >>>>>>>>>> figs/fig-percolation.tex >>>>>>>>>>
% 图 14-1  渗流阈值：团块、跨越簇与序参量
% 连通性按四邻接（von Neumann）核验：
%   (a) 15/36=0.417 格，最大团块 5 格，无左右/上下跨越簇
%   (b) 24/36=0.667 格，唯一左右跨越簇 12 格，其余最大 4 格，无上下跨越
\begin{tikzpicture}[x=1cm, y=1cm, font=\small,
  pt/.style={anchor=center, font=\scriptsize\heiti},
  tick/.style={anchor=north, font=\scriptsize, text=slate},
  leg/.style={anchor=west, font=\scriptsize\kaishu, text=slate},
  ax/.style={-{Stealth[length=2.0mm]}, ink, line width=0.6pt}
]
% 下缘原先 -6.12，(a)/(b)/(c) 三组小标题距页边只剩 1.6--2.5pt。放宽到 -6.35。
\useasboundingbox (0.70,-0.92) rectangle (14.02,-6.35);

\def\s{0.50}
\def\ayl{-5.10}

% ================= 图例 =================
\fill[slate!32] (3.02,-1.42) rectangle (3.32,-1.12);
\node[leg] at (3.40,-1.27) {占据格};
\fill[rust!78] (4.42,-1.42) rectangle (4.72,-1.12);
\node[leg] at (4.80,-1.27) {最大团块};
\draw[slate!45, line width=0.4pt] (5.98,-1.42) rectangle (6.28,-1.12);
\node[leg] at (6.36,-1.27) {空格};
\node[leg] at (7.48,-1.27) {（四邻接）};

% ================= (a) p < p_c =================
\begin{scope}[shift={(0.80,\ayl)}]
  \foreach \i in {0,...,5} \foreach \j in {0,...,5}
    \draw[slate!28, line width=0.25pt] ({\i*\s},{\j*\s}) rectangle ({\i*\s+\s},{\j*\s+\s});
  % 占据格（未高亮）
  \foreach \i/\j in {4/0,2/1,4/1,4/2,2/3,3/3,4/4,5/4,0/5,5/5}
    \fill[slate!32] ({\i*\s},{\j*\s}) rectangle ({\i*\s+\s},{\j*\s+\s});
  % 最大团块：5 格
  \foreach \i/\j in {0/0,1/0,0/1,0/2,1/2}
    \fill[rust!78] ({\i*\s},{\j*\s}) rectangle ({\i*\s+\s},{\j*\s+\s});
  \draw[slate!45, line width=0.5pt] (0,0) rectangle ({6*\s},{6*\s});
\end{scope}
\node[pt, text=ink] at (2.30,-1.72) {$p\approx0.42<p_c$};
\node[pt, text=rust!85] at (2.30,-5.30) {$s_{\max}=5$};

% ================= (b) p > p_c =================
\begin{scope}[shift={(5.20,\ayl)}]
  \foreach \i in {0,...,5} \foreach \j in {0,...,5}
    \draw[slate!28, line width=0.25pt] ({\i*\s},{\j*\s}) rectangle ({\i*\s+\s},{\j*\s+\s});
  % 其余占据格
  \foreach \i/\j in {0/3,2/3,4/3,5/3,0/4,2/4,3/4,5/4,0/5,1/5,3/5,5/5}
    \fill[slate!32] ({\i*\s},{\j*\s}) rectangle ({\i*\s+\s},{\j*\s+\s});
  % 左右跨越簇：12 格
  \foreach \i/\j in {0/0,1/0,2/0,4/0,5/0,0/1,1/1,2/1,3/1,4/1,1/2,3/2}
    \fill[rust!78] ({\i*\s},{\j*\s}) rectangle ({\i*\s+\s},{\j*\s+\s});
  \draw[slate!45, line width=0.5pt] (0,0) rectangle ({6*\s},{6*\s});
  % 两端贯通
  \draw[-{Stealth[length=1.8mm]}, ink, line width=0.6pt] (-0.58,{0.5*\s}) -- (-0.06,{0.5*\s});
  \draw[-{Stealth[length=1.8mm]}, ink, line width=0.6pt] ({6*\s+0.58},{0.5*\s}) -- ({6*\s+0.06},{0.5*\s});
\end{scope}
\node[pt, text=ink] at (6.70,-1.72) {$p\approx0.67>p_c$};
\node[pt, text=rust!85] at (6.70,-5.30) {$s_{\max}=12$};

% p 增大
\draw[ax] (4.10,-3.45) -- (4.90,-3.45);
\node[anchor=north, font=\scriptsize\kaishu, text=slate] at (4.50,-3.58) {$p$ 增大};

% ================= (c) 序参量曲线 =================
\def\ox{8.90} \def\oy{-5.10} \def\ow{4.70} \def\oh{3.00}
\draw[ax] (\ox,\oy) -- ({\ox+\ow+0.15},\oy)
      node[anchor=west, font=\scriptsize, text=ink] {$p$};
\draw[ax] (\ox,\oy) -- (\ox,{\oy+\oh+0.12})
      node[anchor=south, font=\scriptsize, text=ink] {$P_\infty$};
\foreach \v in {0.5,0.6,0.7,0.8} {
  \pgfmathsetmacro\xx{\ox+\ow*(\v-0.5)/0.30}
  \draw[slate!60, line width=0.4pt] (\xx,\oy) -- (\xx,\oy-0.08);
  \node[tick] at (\xx,\oy-0.10) {\v};
}
\foreach \v in {0,1} {
  \pgfmathsetmacro\yy{\oy+\oh*\v}
  \draw[slate!60, line width=0.4pt] (\ox,\yy) -- (\ox-0.08,\yy);
  \node[tick, anchor=east] at (\ox-0.10,\yy) {\v};
}
% 临界点
\pgfmathsetmacro\pcx{\ox+\ow*(0.5927-0.5)/0.30}
\draw[slate!60, line width=0.5pt, dashed] (\pcx,\oy) -- (\pcx,{\oy+\oh-0.05});
% 主曲线：p<p_c 恒为 0，过 p_c 按 β=5/36 上升
\draw[rust!85, line width=1.3pt] (\ox,\oy) -- (\pcx,\oy);
\draw[rust!85, line width=1.3pt]
      plot[domain=0.5927:0.80, samples=90]
      ({\ox+\ow*(\x-0.5)/0.30}, {\oy+\oh*((\x-0.5927)/0.4073)^0.13889});
\fill[rust!85] (\pcx,\oy) circle (1.5pt);
\node[anchor=south west, font=\scriptsize, text=ink] at ({\pcx+0.16},{\oy+\oh+0.02})
      {$p_c\approx0.5927$};
\node[anchor=west, font=\scriptsize, text=rust!85] at (10.60,-3.90)
      {$P_\infty\propto(p-p_c)^{5/36}$};

% ================= 栏标题（一律在内容下方）=================
\node[anchor=north, font=\footnotesize\heiti, text=slate] at (2.30,-5.72) {(a) 孤立团块};
\node[anchor=north, font=\footnotesize\heiti, text=slate] at (6.70,-5.72) {(b) 跨越簇};
\node[anchor=north, font=\footnotesize\heiti, text=slate] at (11.25,-5.72) {(c) 序参量};
\end{tikzpicture}

% <<<<<<<<<< figs/fig-percolation.tex <<<<<<<<<<

\caption{渗流阈值附近的相变。占据比例低于阈值时，系统里只有彼此分离的小团块，扰动无法跨越；越过阈值后出现贯通整体的跨越簇，行为性质随之改变。右下角的序参量 $P_\infty$ 在阈值附近迅速上升，过渡区的宽窄由系统规模决定。}
\label{fig:percolation}
\end{figure}

这件事在数学模型里可以很干净，在现实系统中却总是混杂发生。真实组织不是随机图，城市交通不是静态网络，供应链不是只由节点和边组成，金融系统也不只是交易关系的集合。现实系统中还有时间延迟、激励结构、权力分布、标准接口、责任边界、信息不对称和人的预期。正因为如此，现实中的凝胶化常常更难被看见：它不是在某一天明确跨过一个可测阈值，而是在连接、标准、流程、平台、指标、权限和自动化不断叠加之后，突然让人发现这个系统已经不能局部改了。低耦合状态下，一个变化可以被隔离，一个失败可以被局部修复，一个新技术可以小范围试验，系统有空隙，有缓冲，有局部自由度；高耦合状态下，连接贯通之后，局部变化会进入全局通道，一个接口改变，会影响上下游系统，一个指标调整，会改变多个主体的行为，一个模型接入，会牵动数据、审核、执行、合规和责任链条。系统不再允许局部随意流动，而表现出整体性的硬化。这种硬化不是来自某个节点本身，而是来自贯通结构：一张松散的网可以拉动其中一小段，一张绷紧的网则不同，一个点被拉动，整张网都会响应。复杂系统一旦越过某种连接阈值，局部自由度就会被整体结构接管，一个部门想快，必须让其他部门吸收，一个模块想升级，必须让接口、数据、权限、培训和责任一起更新，一个技术想部署，必须同时改变流程、标准、激励和风险分配。

所以，连接越多，并不必然意味着系统越灵活。连接提高协同，也提高牵连；连接提高效率，也提高传播风险；连接减少局部摩擦，也可能增加系统摩擦。一个系统越能无缝协同，越可能在出错时无缝传播冲击。流动性和脆弱性，常常来自同一张连接网络。这个判断反直觉，因为现代化通常把连接看成进步：信息系统要连接，供应链要连接，城市要连接，金融要连接，组织要打通数据孤岛，AI要接入流程、接口、工具和执行系统。连接带来效率，这是事实，没有连接，系统只是分散的部件，无法协同，无法规模化，无法产生复杂功能。但连接也会改变失败的性质：低连接系统的失败常常是局部坏掉，高连接系统的失败可能是级联、放大、同步和全局停摆。系统不是因为连接而天然更强，而是因为连接而获得了更强的传播能力，传播能力既能传递协同，也能传递错误。

相互依赖网络研究把这个问题推得更清楚。Buldyrev等人在关于相互依赖网络的研究中指出，一个网络中的节点失效可能引发另一个网络中依赖节点失效，这种失效可以递归传播，导致级联；在某些条件下，去除很小一部分节点就可能造成相互依赖网络的巨大破碎\cite{buldyrev2010}。这不是要说所有现实系统都会这样崩溃，而是说明一个重要机制：当网络不是单层的，而是由电力、通信、金融、物流、组织和信息系统彼此支撑时，脆弱性不只来自每一层内部，还来自层与层之间的依赖关系。独立网络中可以承受的冲击，进入相互依赖网络后，可能变成跨层级联。层间依赖的耦合强度决定了级联的规模，如果一层对另一层的依赖是单向的，冲击的传播方向也随之固定，这让级联的路径可以被预测，也让薄弱环节可以被定位。

\subhead{级联的实例与代价}

复杂系统中的风险常常被低估。评估习惯是分别检查节点是否可靠、模块是否稳定、流程是否合规、模型是否准确，却容易忽略它们连接后的整体相态。一个模块在局部测试中可靠，不等于它在全系统中可靠；一个模型在离线指标中准确，不等于它在组织流程中可靠；一个策略在单点实验中有效，不等于它在多人响应、反馈延迟和激励改变之后仍然有效。凝胶化系统的风险不在部件单独坏，而在部件之间的牵连方式改变了错误传播。

大型组织最容易显示这种凝胶化。一个小团队更换工具，可能只需要几次沟通；一个大型组织更换同样工具，却会牵动权限、数据、安全、预算、培训、审计、绩效和合规。工具本身没有变复杂，但它进入的网络已经不同。低耦合环境中的工具选择，到了高耦合组织中，会变成结构迁移。很多组织数字化失败，不是因为软件不好，而是因为软件触动了已经凝胶化的流程网络：每个节点都同意提高效率，但没有哪个节点愿意单独承担接口变化、责任重分配和旧流程失效带来的风险。

城市交通也是如此。一个路口的信号灯优化，在低密度区域可能直接改善通行；在高密度城市网络中，一个路口变快，可能只是把车辆更快送到下一个堵点，局部速度提高，未必带来全局速度提高。第一章讲过的布雷斯悖论是这件事的极端形式：增加一条道路，在某些网络条件下反而使整体交通变差，而原因不在道路本身，在于个体路径选择和网络均衡会重新分配流量。交通不是管道，而是带有策略主体的网络，一个连接的增加，会改变许多人的行为：连接改变结构，结构改变行为，行为再反过来改变系统状态。

供应链中的凝胶化更直接。在低耦合供应链中，一个零件延迟可能只是局部延迟；在高度精益、低库存、全球分工和标准化接口密集的供应链中，一个小零件短缺就可能让整条生产线停摆。系统在正常状态下非常高效，因为冗余被压缩，连接被强化，库存被降低；但也正因为如此，局部冲击更容易贯通全局。效率和脆弱性不是相反物，而可能是同一个结构的两面。

苏伊士运河的堵塞就是一个直观案例。2021年Ever Given搁浅使苏伊士运河受阻近一周，后续共有422艘等待通行的船只通过后才清空积压\cite{suez2021}。这条航道平时显得像全球贸易的效率通道，因为它缩短了亚欧航线，减少时间和燃料；但当它被一艘船阻断时，全球物流就突然显露出高耦合网络的另一面。新西兰外交贸易部的报告称，约12\%的全球贸易和约30\%的全球集装箱运输量经过苏伊士运河\cite{nzmfat2021}。这个案例的意义不在于夸大一次事故，而在于说明：高效通道本身就是结构性依赖，平时它减少摩擦，危机时它放大堵塞。

电力系统提供了更严肃的级联图像。2003年8月14日，美国东北部和加拿大安大略发生大停电，官方报告称影响区域估计涉及5000万人和61800兆瓦电力负荷，部分地区停电持续到4天\cite{uscanada2004}。更关键的是，这不是单一设备坏掉造成的简单事故，而是保护、调度、信息显示、植被管理、系统监控和区域协调共同失效后的级联事件。报告还指出，在某个关键线路跳闸之后，级联进入多个州和省份，后续扩散在数分钟内完成，其中大部分事件发生在级联最后的极短时间窗口内\cite{uscanada2004}。电网平时作为高度互联结构提高供电可靠性和资源调度能力，但当系统进入不安全状态，互联也会成为故障扩散通道。

金融系统同样如此。金融连接让资本流动更快，资产定价更连续，风险转移更灵活，但当银行、基金、保险、衍生品、抵押品、评级、算法交易和市场预期彼此绑定，一个局部价格下跌就可能触发保证金追加、强制平仓、流动性收缩、恐慌抛售和监管反应。平稳时期，连接让系统像液体一样流动；危机时期，同一张连接网络会让系统像凝胶一样整体收缩。2010年5月6日，美国股市在几分钟内出现盘中暴跌，道琼斯指数一度下跌近千点，随后又迅速回弹，美国证券交易委员会与商品期货交易委员会在联合报告中指出，自动交易系统在流动性突然枯竭时执行的大额卖出指令，与高频做市商的快速退出共同放大了这次波动\cite{sec2010findings}。价格在几分钟内被推离基本面又被拉回，损失与收益却已经在真实账户中完成结算，这正是一个局部指令在高耦合网络中变成系统事件的形状。

Knight Capital的交易事故显示了自动化系统在高耦合金融环境中的速度：美国证券交易委员会披露，2012年8月1日，Knight Capital的自动化交易路由器在开盘后45分钟内向市场发送了超过400万个错误订单，试图执行仅212个客户订单，交易超过3.97亿股，获得数十亿美元非预期头寸，并最终造成超过4.6亿美元损失\cite{sec2013knight}。这不是一般意义上的“一个程序出错”，它说明在高频、自动化、市场接入和实时执行耦合在一起的系统中，软件缺陷可以迅速转化为市场行为、资产头寸和公司生存危机。比特错误一旦接入金融执行网络，就不再只是比特错误。

供应链的级联在近年的芯片短缺中又一次显形。汽车行业长期按零库存方式组织生产，芯片以长周期、低价值的形式分散在几十个控制单元里，当晶圆产能被其他用途挤占之后，缺少的往往不是最贵的那颗芯片，而是几美元的车身控制模块，缺少一个模块足以让整车无法下线。产能恢复本身需要时间，晶圆厂的扩产以年计，而需求端的重启要快得多，两端的时间尺度差把短缺拉长到了几个季度。这个案例把凝胶化的两个特征同时展示出来：断点在价值最低的位置，恢复时间由最慢的那一层决定。

教育系统也有这种结构，只是它的级联不表现为停摆，而表现为行为重排。一个评价指标改变，可能不只是改变学校内部管理，还会改变教师教学、学生策略、家长投入、培训机构、学校排名和地方政策。指标一旦进入高耦合网络，就不再是测量工具，而成为行为牵引器，它改变的是一整套行为场。很多改革在局部看是合理的：更清晰的指标、更透明的排名、更精细的评价、更频繁的数据反馈。但如果这些指标进入一个由升学、资源分配、声誉竞争和家长期待构成的凝胶化系统，它就会改变系统的策略结构，测量越密，行为越围绕测量组织，评价越频繁，真实目标越容易被指标替代\cite{goodhart1975problems,campbell1979assessing}。

\subhead{AI的三种局部结果}

AI进入这些系统时，面对的不是单个任务，而是已经接近凝胶化的网络。AI在低耦合场景中表现轻盈：写文案、生成摘要、辅助检索、补全代码、制作草图，这些任务失败可隔离，输出可回滚，责任范围有限。但当AI进入医疗、金融、交通、教育、能源和公共治理，它的输出会立刻接入数据、流程、审核、执行、责任和人的行为，它不再只是一个工具，而成为网络中的扰动源。这也是为什么大模型在比特空间中显得异常快，而进入现实系统后速度会突然下降：比特任务往往允许高频试错，错了可以重写，结果可以撤回，版本可以更新，候选可以并行生成；高耦合系统则不同，医疗建议会影响检查、治疗、保险和法律责任，交通调度会影响车辆、道路、乘客、事故和公共秩序，金融决策会影响资产、风险敞口、监管和市场预期，教育评价会影响家庭投入、学校行为和学生的路径选择。AI输出一旦接入这些系统，就从“候选生成”变成“结构扰动”。这时，真正的瓶颈不是模型能否生成，而是系统能否吸收。

在凝胶化系统中，AI的局部优化可能出现三种结果。第一，它被吸收：系统原有结构不改变，AI只是成为边缘辅助工具，提高局部效率，但不改变整体状态，比如自动生成会议纪要、辅助起草邮件、初步分类材料，这些应用可以被原有流程吸收，它们改善局部体验，却不触动责任结构、权力结构和行动链条。第二，它被抵消：一个环节因为AI变快，另一个环节成为瓶颈。报告生成更快，审核更慢；方案更多，决策更难；预测更频繁，执行能力跟不上；客服回复更快，复杂投诉升级更多；代码候选更多，测试、集成和安全审查压力更大。系统总速度没有提高，只是堵点换了位置，低耦合任务中局部加速就近似等于收益，凝胶化系统中局部加速常常变成队列迁移。第三，它被放大为新问题：自动化提高局部效率，却降低异常恢复能力；推荐系统提高点击率，却放大情绪传播；高频算法提高交易速度，却增加市场震荡；教育评价系统提高测量频率，却放大应试行为；医疗分诊模型提高处理速度，却可能把偏差嵌入更大流程。AI不只是替换某个动作，而可能改变系统响应方式，它让某种行为更便宜，也让这种行为更容易规模化和制度化。

三种结果的分界不在模型能力上，而在系统是否保留了吸收新输入的余地。同样的模型放进两种环境，结果可以完全相反：在责任可以隔离、失败可以回滚的位置，它表现为效率提升；在责任无法隔离、失败无法回滚的位置，它表现为风险转移。判断落在哪一边，要看改动是否触及节拍源、是否跨过组织边界、是否让某一层的输入频率超过了它的更新时间。这三条检查做完，一项应用会被吸收还是会被放大，通常已经可以判断。

因此，在凝胶化系统中，技术评估不能只看局部性能，真正的问题是系统如何吸收这个性能。一个工具进入低耦合场景，像水流进入河道；进入凝胶化系统，则像外力作用在一张绷紧的网，输入相同，响应完全不同。模型准确率、任务完成率、响应速度、用户满意度，都只是局部指标。系统级问题是：错误如何传播？异常如何停止？责任如何回流？瓶颈是否迁移？人是否改变策略？旧流程是否被迫重排？系统有没有足够的缓冲和隔离，让试错不直接贯通全局？

\subhead{规模化与近可分解}

很多技术试点容易、规模化困难。试点通常发生在局部团簇中，系统连接被人为切断，失败影响被隔离，参与者被选择，资源被特别配置，组织注意力高度集中。可一旦扩展到全系统，连接重新接通，隐藏耦合就会出现：数据标准不同，流程不同，责任边界不同，人的接受度不同，监管要求不同，旧系统接口不同。规模化不是把试点复制很多次，而是让试点穿过一个更大的网络。很多技术不是在功能上失败，而是在贯通网络时变硬。

企业数字化、智慧城市、智能医疗和AI政务都常遇到这个问题。一个部门试点顺利，不代表跨部门运行顺利；一个医院科室验证有效，不代表全院流程能接住；一个城市路段算法优化，不代表全城交通改善；一个审批环节自动化，不代表整个公共服务链条变快。试点把系统切薄了，规模化把系统重新接厚了。凝胶化系统的阻力，往往就在这个“重新接厚”的过程中显影。

这里要特别警惕一种工程幻想：只要连接更多数据、打通更多接口、接入更多模型，系统就会自动变聪明。现实恰好相反。连接越多，系统可能越能感知，也可能越难改变；越能自动响应，也可能越难停下；越能全局协同，也可能越容易全局共振。一个高度互联的系统，不一定更有智慧，它只是更少隔离，它可以更快发现问题，也可以更快传播错误，可以更快优化，也可以更快失稳。

凝胶化系统提醒我们，更多连接并不总是更好。连接过少，系统无法协同；连接过密，系统失去弹性。好的结构不是连接最多的结构，而是在连通性与隔离性之间保持张力：它既允许信息、资源和行动传递，也允许局部错误被限制，局部试验不牵动全局，失败有缓冲带。这正是模块化、冗余、沙盒、分层接口、灰度部署和断路器的价值，它们的共同作用不是追求更多连接，而是控制连接的贯通方式，让必要的协同发生，同时让试验、错误和异常不至于直接穿透全系统。它们不是保守设计，而是复杂系统保持可塑性的条件，一个没有隔离结构的系统，短期看顺滑，长期看脆弱；一个保留冗余和模块边界的系统，短期看低效，长期看反而更能吸收变化。

这种思想和Herbert Simon所说的\term{近可分解}系统相通。复杂系统之所以能够演化和管理，往往不是因为所有部分完全打通，而是因为系统存在层级、模块和相对独立的子结构\cite{simon1962}。子系统之间有连接，但连接不是无边界贯通，局部变化可以先在子系统内部被吸收、测试和修正，再通过接口影响更大系统。没有这种近可分解性，系统每一次局部调整都会变成整体重排，演化就会极端昂贵。AI进入现实系统，也必须遵守这种结构逻辑：低风险比特任务可以快速试错，高耦合系统需要分层验证、局部隔离、失败回滚和责任缓冲。医疗、交通、金融、电网、教育和公共治理需要的不只是更强模型，还需要防止模型输出直接贯通全系统的结构设计，否则，比特层的高频输入会冲击低频结构，造成堆积、震荡和系统性风险。

\subhead{介质梯度与技术黏滞}

这里与技术黏滞的主线直接相接。技术黏滞不是系统不愿意进步，也不是人类保守，而是智能输出进入现实系统后，必须面对资源、时间和结构耦合的阻尼。凝胶化系统显示的是结构维度的非线性：连接一旦贯通，局部能力不再按局部收益结算，而要按系统后果结算，一个AI模型在单任务上提升20\%，并不意味着系统效率提升20\%。它可能被吸收，可能被抵消，也可能触发新的协调、验证和责任成本。

从比特、原子和生命三个介质看，凝胶化程度也不同。比特空间最容易快速连接，接口、协议、API、数据库和模型调用可以迅速形成高频网络，因此它也最容易发生信息级联、推荐共振、舆论放大和软件级事故。原子空间连接较慢，但一旦进入交通、电网、制造、物流和城市基础设施，错误更难回滚，恢复需要材料、能量、维修和时间。生命介质更重，因为生命过程具有发育、代谢、免疫、繁殖和生态耦合，一个干预进入生命网络，可能不是单点效果，而是长期路径改变。介质越重，凝胶化后的代价越难撤回。

复杂网络不是平滑可改的。它可能在低耦合区柔软，在高耦合区突然变硬。技术进入这样的系统时，不能只问局部能力是否足够，还要问系统是否已经越过某种连接阈值：局部变化还能否被局部吸收？错误还能否被隔离？新方案还能否渐进接入？反馈还能否稳定返回？如果这些问题没有回答，所谓部署只是把局部能力推入一张已经绷紧的网。高耦合系统的阻力不是逐渐增加的，而可能在某个临界区域之后突然变硬，连接一旦贯通，局部变化就不再是局部变化。

因此，技术进入高耦合系统时，真正需要处理的并不只是阻力大小，而是连接结构本身。连接一旦贯通，局部变化就会获得系统后果，系统也会从可局部调整，转向整体性硬化。智能时代最危险的误判，是把局部能力当成系统能力，把连接当成灵活，把自动化当成可控，把试点成功当成规模化成功。凝胶化系统要求我们反过来思考：不是先问技术能做什么，而是先问系统是否还有足够的空隙、缓冲和边界，允许这个技术以可吸收、可回滚、可承责的方式进入。

系统一旦越过阈值，它的状态就不再只由当前连接决定。此前已经投入的资产、已经采用的接口、已经形成的习惯和已经被改变的行为，都留在结构里，并不会随着目标的更换而消失。为什么这些残余不能在技术撤回时一起清除，决定了一个高耦合系统还有没有恢复的余地。

\subsection{迟滞效应与结构记忆}

系统变硬之后，局部调整的代价上升，但变化并没有停止。已经完成的投资、已经采用的接口、已经形成的习惯，仍然在系统内部起作用，它们不会因为当前目标改变而消失。现实系统有记忆。这不是修辞，而是系统演化中的真实机制。一个系统不是每一刻都从零开始选择最优状态，它的当前形态，是过去投资、过去制度、过去接口、过去错误、过去妥协和过去习惯沉积下来的结果。道路、医院、学校、企业、软件、产业链、法律制度和城市空间，都不是抽象设计的结果，而是历史路径逐层叠加的结果。

\subhead{物理残余}

但结构记忆不是“历史包袱”这么简单。\term{结构记忆}是过去在系统中留下的物理残余：它改变了路径代价，改变了接口形态，改变了连接强度，改变了尺度传导，改变了回退方向，改变了系统后来能够如何被改变。过去之所以没有消失，不是因为系统记得它，而是因为系统已经被它改过了。一条道路修成之后，留下的不只是道路本身，它会改变土地价值、居住分布、商业位置、通勤习惯和后续基础设施；一个软件接口被采用之后，留下的不只是技术选择，它会改变数据格式、开发流程、兼容要求和后续系统的设计边界；一个评价指标被长期使用之后，留下的不只是一个数字，它会改变人的努力方向、组织的资源分配和后来评价新方案的标准。这些东西共同说明：时间会把选择变成结构，把结构变成惯性。

“物理残余”这个说法在这里是字面意义上的。铁路的轨距一旦铺定，车辆、站台、隧道断面、维修设备和调度规则都围绕它定型，之后的变更成本不再由钢轨本身决定，而由整套附属设施决定。1435 毫米的标准轨之所以成为标准，是历史竞争的结果，而不是几何上的必然；采用 1520 毫米的独联体国家、采用 1676 毫米的印度和采用 1668 毫米的西班牙，都保留了自己的制式，并在边境上依靠变轨装置或换装作业来对接。制式之间的差别只有几百毫米，转换代价却以十亿计。电网的电压等级与频率、家用电器的额定电压、螺栓的螺纹标准、集装箱的尺寸，都有同样的性质：它们最初可能是任意选定的，一旦被大规模生产、认证和培训体系采纳，就从选择变成了环境。

\subhead{路径依赖与迟滞}

惯性的第一层结果，是前进方向被已经走过的路约束。这也是\term{路径依赖}的基本含义。现实系统不是沿着抽象最优路径运动，而是在已经形成的路径上继续前进，过去的选择改变了未来可选项\cite{david1985clio}。一个系统曾经为了某个目标建立了标准、接口、资产和习惯，这些东西后来就不再只是工具，而变成未来选择的边界，它们会让某些改变更容易，让另一些改变更昂贵。所以，路径依赖不是“人们习惯走老路”。它的第一性含义是：走过的路会改变后来的路。系统从A走向B，并不是把自身从一个点搬到另一个点，这个过程会改变系统本身：人员被重新训练，资产被重新配置，接口被重新设计，数据被重新格式化，责任被重新划分，利益被重新分配。系统进入B之后，原来的A也不再存在，此时如果想回到过去状态，往往不能沿原路返回。

这就是\term{迟滞}。迟滞意味着，变化路径具有方向性：从A到B的代价，不等于从B回到A的代价。一个企业迁移到新的软件系统，可能花费多年清理数据、重建流程、培训员工、适配接口；一旦迁移完成，再回到旧系统并不是简单卸载新系统，因为旧人员、旧流程、旧数据习惯和旧供应商关系已经改变。一个城市形成某种道路和住宅格局之后，即使后来发现更优规划，也不能像修改图纸一样回到早期状态，已有建筑、土地权属、居民生活和基础设施，会把过去选择固化在空间中。不可逆并不总是绝对不能回退，而是回退本身已经成为新的工程。系统走过一条路之后，路会改变系统，回退不是倒放，而是在被改变后的系统里重新改变：名义上回到A，结构上已经不是原来的A，材料、接口、习惯、信任和责任都已经留下残余。

迟滞在物理测量里有严格的形式，这有助于看清它不是拖延。铁磁材料在外加磁场下磁化，撤去外场后磁化强度并不回到零，保留下来的部分称为剩磁；要让磁化归零，必须施加一个反向磁场，其大小称为矫顽力。把磁化强度随外场变化的曲线画出来，得到一个闭合的回线，回线的面积对应每循环一次所耗散的能量，这就是变压器和电机铁芯发热的来源。整个过程里，材料当前的状态不是外场当前值的函数，而是外场走过的整段历史的函数，同一条外场路径走回去，走的是另一条曲线。任何具有这种回线特征的系统，其状态都不能由当前输入单独读出，撤去输入也不意味着回到原点。

\begin{figure}[ht]
\centering
% >>>>>>>>>> figs/fig-hysteresis.tex >>>>>>>>>>
% 图 15-1  迟滞：状态由路径而非瞬时输入决定
% (a) X-Y 回线：下支 Y=0.02+0.10X，上支 Y=0.80+0.18X，双稳区 X 在 [0.35,0.62]
%     升程在 X_up=0.62 上跳，降程在 X_down=0.35 下跳 -> 回线闭合
% (b) X(t) 三角波与 Y(t) 阶跃，跳变时刻 t1=3.72、t2=9.90（与 0.62/0.35 一致）
\begin{tikzpicture}[x=1cm, y=1cm, font=\small,
  ax/.style={-{Stealth[length=2.0mm]}, ink, line width=0.6pt},
  tick/.style={font=\scriptsize, text=slate}
]
% 左边与下边原来各切掉约 0.6pt 与 2.5pt：纵轴刻度字贴左缘，
% 底部 (a)/(b) 两个小标题贴下缘。两侧各放 0.16cm。（用 _figclip.py 量出）
\useasboundingbox (0.52,-0.70) rectangle (14.06,-6.74);

\def\ax{0.95} \def\ay{-5.60} \def\aw{5.30} \def\ah{4.30}
\def\bx{7.85} \def\by{-5.60} \def\bw{5.30} \def\bh{4.30}

% 支路参数（屏幕坐标）
\def\xdn{2.805}          % X_down = 0.35
\def\xup{4.236}          % X_up   = 0.62

% ================= (a) X-Y 回线 =================
\draw[ax] (\ax,\ay) -- ({\ax+\aw+0.35},\ay) node[anchor=west, font=\scriptsize, text=ink] {$X$};
\draw[ax] (\ax,\ay) -- (\ax,{\ay+\ah+0.15}) node[anchor=south, font=\scriptsize, text=ink] {$Y$};

% 回线内部（双稳区）先铺底色
\fill[rust!10]
  (\xdn,-5.3635) -- (\xup,-5.2474) -- (\xup,-1.6801) -- (\xdn,-1.8891) -- cycle;
% 下支（升程走的一支）
\draw[deepblue!85, line width=1.3pt] (\ax,-5.514) -- (6.25,-5.084);
% 上支（降程走的一支）
\draw[terra!90, line width=1.3pt] (\ax,-2.160) -- (6.25,-1.386);
% 跳变箭头
\draw[-{Stealth[length=2.2mm]}, rust!90, line width=1.0pt] (\xup,-5.2474) -- (\xup,-1.6801);
\draw[-{Stealth[length=2.2mm]}, rust!90, line width=1.0pt] (\xdn,-1.8891) -- (\xdn,-5.3635);
% 支路上的方向箭头
\draw[-{Stealth[length=1.8mm]}, deepblue!85, line width=0.8pt] (1.48,-5.471) -- (2.01,-5.428);
\draw[-{Stealth[length=1.8mm]}, terra!90, line width=0.8pt] (5.72,-1.463) -- (5.08,-1.556);
% 双稳区标注（竖排在回线内部）
\node[font=\scriptsize\kaishu, text=rust!90, rotate=90] at (3.52,-3.52) {双稳区};

% 坐标刻度
\foreach \v/\lb in {0/0, 0.35/{$X_\downarrow$}, 0.62/{$X_\uparrow$}, 1/1} {
  \pgfmathsetmacro\xx{\ax+\aw*\v}
  \draw[slate!55, line width=0.4pt] (\xx,\ay) -- (\xx,\ay-0.08);
  \node[tick, anchor=north] at (\xx,\ay-0.10) {\lb};
}
\foreach \v in {0,1} {
  \pgfmathsetmacro\yy{\ay+\ah*\v}
  \draw[slate!55, line width=0.4pt] (\ax,\yy) -- (\ax-0.08,\yy);
  \node[tick, anchor=east] at (\ax-0.10,\yy) {\v};
}

% ================= (b) 驱动与响应的时间历程 =================
\draw[ax] (\bx,\by) -- ({\bx+\bw+0.60},\by) node[anchor=west, font=\scriptsize, text=ink] {$t$};
\draw[ax] (\bx,\by) -- (\bx,{\by+\bh+0.15}) node[anchor=south, font=\scriptsize, text=ink] {$X,\,Y$};
\foreach \v/\lb in {0/0, 6/6, 12/12} {
  \pgfmathsetmacro\xx{\bx+\bw*\v/12}
  \draw[slate!55, line width=0.4pt] (\xx,\by) -- (\xx,\by-0.08);
  \node[tick, anchor=north] at (\xx,\by-0.10) {\lb};
}
\foreach \v in {0,1} {
  \pgfmathsetmacro\yy{\by+\bh*\v}
  \draw[slate!55, line width=0.4pt] (\bx,\yy) -- (\bx-0.08,\yy);
  \node[tick, anchor=east] at (\bx-0.10,\yy) {\v};
}
% X(t)：三角波
\draw[slate!75, line width=1.0pt, dash pattern=on 3pt off 1.6pt]
      (\bx,\by) -- ({10.50},{\by+\bh}) -- ({\bx+\bw},\by);
% Y(t)：在 t1 上跳、t2 下跳的阶跃
\def\tx{9.493}
\def\ty{12.223}
\draw[rust!90, line width=1.4pt]
      (\bx,-5.342) -- (\tx,-5.342) -- (\tx,-1.644) -- (\ty,-1.644) -- (\ty,-5.342) -- ({\bx+\bw},-5.342);
% 跳变时刻
\draw[slate!45, line width=0.4pt, dashed] (\tx,\by) -- (\tx,{\by+\bh});
\draw[slate!45, line width=0.4pt, dashed] (\ty,\by) -- (\ty,{\by+\bh});
\node[font=\scriptsize\heiti, text=rust!90, anchor=north] at (\tx,-5.72) {$t_1$};
\node[font=\scriptsize\heiti, text=rust!90, anchor=north] at (\ty,-5.72) {$t_2$};
% 图例（图内左上空白）
\draw[slate!75, line width=1.0pt, dash pattern=on 3pt off 1.6pt] (7.95,-1.66) -- (8.40,-1.66);
\node[font=\scriptsize\kaishu, text=slate, anchor=west] at (8.47,-1.66) {驱动$X$};
\draw[rust!90, line width=1.4pt] (7.95,-2.06) -- (8.40,-2.06);
\node[font=\scriptsize\kaishu, text=slate, anchor=west] at (8.47,-2.06) {响应$Y$};

% ================= 栏标题（一律在内容下方）=================
\node[anchor=north, font=\footnotesize\heiti, text=slate] at (3.52,-6.22) {(a) 迟滞回线};
\node[anchor=north, font=\footnotesize\heiti, text=slate] at (10.50,-6.22) {(b) 时间历程};
\end{tikzpicture}

% <<<<<<<<<< figs/fig-hysteresis.tex <<<<<<<<<<

\caption{迟滞：状态由路径决定，而不由瞬时输入决定。左侧的输入-状态回线上，同一个输入值对应两个不同的状态；右侧把输入与响应的变化放在同一条时间轴上，撤去输入之后状态停在新的位置，要回到起点必须付出反向的代价。}
\label{fig:hysteresis}
\end{figure}

\subhead{冻结的自由度}

这使结构记忆具有物理意味。材料经过反复受力，会留下疲劳和残余应力；磁性材料经过外场，会出现磁滞；玻璃冷却时，会冻结部分无序结构；复杂网络经过长期使用，会形成高权重边、关键接口和低阻路径。不同系统的材料不同，但机制相似：过去改变了内部自由度的组织方式，使系统不再像初始状态那样可任意重排。\key{结构记忆，就是冻结下来的自由度。}一个系统一开始可能有许多可能连接、许多可能标准、许多可能路径，但一旦某些连接被反复使用，某些接口被大量兼容，某些标准被广泛采用，某些资产被长期投入，系统的自由度就会减少。不是所有方向都还同样可行，不是所有路径都还同样便宜，历史把一部分可能性转化为低阻路径，也把另一部分可能性推成高阻边界。

玻璃的冷却过程给出了这个机制最干净的版本。液态时原子可以相对自由地重新排列，体系有大量可能的构型；温度下降时，重新排列所需的时间按指数增长，到某个温度以下，原子已经来不及在可观测的时间内完成重排，构型就被冻结在当时的无序状态里。若冷却得慢，体系能落到更低的能量状态；若冷却得快，则停在更高的能量状态，形成玻璃。两者成分完全相同，性质却不同，差别只在经过的路径。结构记忆的强度因此取决于两件事：当时的自由度有多少，以及变化发生得有多快。

旧系统如此顽固，症结就在这里。旧软件不好用，但所有数据、权限、报表和人员训练都围绕它；旧标准不完美，但设备、供应商、认证和客户关系都围绕它；旧制度有缺陷，但资源分配、晋升路径和风险责任都围绕它。它们不是最优，却已经成为低阻路径；低阻路径长期被使用，就会继续吸引投资、训练、接口和预期，最后把自己加固成结构。局部最优不是点的性质，而是周围能量地形的性质。系统停在一个地方，不一定因为那里最好，而是因为离开它需要同时改变太多耦合：接口要改，标准要改，数据要改，流程要改，习惯要改，责任要改。每一条边都增加一点代价，合起来就形成一个稳定盆地，系统不是不知道别处更好，而是从当前盆地爬出去太重。

\subhead{尺度过滤与结构屏蔽}

结构记忆并不只发生在单一尺度上。复杂系统通常由个体、流程、组织、制度、产业和社会环境组成，一个变化在低层发生，并不意味着它会原样传到高层。上层结构会过滤、吸收、改写或屏蔽下层变化，只让与自身逻辑兼容的部分继续扩散，很多微观变化是真实的，但它们没有成为宏观变量。这里可以借用重整化群的直觉：重整化群在物理学中的一个重要启发是，当我们从微观尺度观察宏观行为时，许多微观细节会被粗粒化掉，只有少数对宏观状态真正起决定作用的变量保留下来。不是每一个微观扰动都能改变宏观相态，大量局部差异会在尺度转换中被平均、过滤或屏蔽\cite{wilson1971renormalization}。

这套方法在临界现象里的用法值得留意。系统处在相变点附近时，不同尺度的涨落呈现相似的结构，把尺度放大若干倍之后，统计特征保持不变，物理量随尺度的变化因此可以被写成幂律，少数几个临界指数就足以刻画整个宏观行为，具体材料的微观细节反而不进入结果。由此可以理解，不同的物理系统在临界点附近会表现出相同的指数：微观差别在尺度变换中被抹掉了。把这个直觉搬到社会技术系统里需要谨慎，因为它不是严格模型，但它提示了一个方向：宏观行为由少数在尺度变换下仍然有效的变量决定，而这些变量往往不是最容易观察到的那些。

社会技术系统中也有类似现象。一个教师用AI改善备课，当然可能提高个人效率，但如果课程标准、考试制度、学校评价、升学竞争和资源分配不变，这种微观改良很难改变教育系统的宏观状态；一个医生用AI提高某类判断效率，也不等于医院整体效率提升，如果科室结构、医保支付、病历流程、法律责任和患者流转机制不变，局部效率会在系统尺度上被吸收；一个车间完成局部自动化，也不等于产业链升级，如果订单结构、供应商能力、物流系统、质量标准和资本投入没有改变，微观优化很难越过宏观瓶颈。重整化群在这里不是严格数学模型，而是一种尺度转换直觉：技术改良要成为系统改变，必须改变那些在宏观尺度上仍然有效的结构变量；否则，局部变化会被系统粗粒化掉，成为个人经验、局部样板或短期项目，而不是整体状态转移。这就是\term{结构屏蔽}。系统不是否定微观创新，而是把微观创新过滤成与自身结构兼容的形状，很多改革之所以“局部有用，整体无感”，并不是局部实践没有价值，而是它没有穿过跨尺度过滤，没有进入评价机制、资源配置、责任边界和制度流程，它没有成为宏观变量。

AI在这里特别容易制造错觉。它让个体能力突然增强：一个人可以更快写作、检索、编程、设计、总结和分析，于是人们容易认为组织、行业和制度也会同步增强。但个体增强不等于结构改变，系统是否改变，取决于这些能力是否被重新接入流程、评价、资源和责任结构；如果没有穿过这些层级，它们就会被宏观结构吸收为局部经验。节点变强，不等于网络改变。一个模型让个人更快完成原有任务，提升的是节点性能；只有当它改变流程接口、责任分布、资源配置和组织边界时，它才改变结构。高耦合系统的瓶颈不一定在节点能力，而在边的组织方式：接口不变，流程不变，权责不变，评价不变，旧结构就会吸收新的局部能力。这也是很多AI应用的真实边界，AI可以制造大量微观扰动，更多方案、更多文本、更多分析、更多自动化脚本、更多局部效率，可是这些扰动能否成为结构变化，取决于它们是否改变了系统的有效耦合；没有改变接口、标准、流程、激励和责任结构的技术，常常只能成为经验、补丁、试点、外挂或报告，而不是系统变革。

判断一项改良有没有穿过尺度过滤，可以看它是否进入了三类变量：评价什么、分配什么、由谁承担。这三类变量在尺度变换下的存活率最高，也最不容易被平均掉。一项技术如果只改变了做事的速度，而没有改变任何一类变量的取值，它在个人层面是真实的效率提升，在系统层面几乎必然被吸收。反过来，一项技术即使幅度不大，只要改变了评价口径或责任归属，就会在更高尺度上留下痕迹，并因此引发一系列后续调整。

\subhead{撤回不等于恢复}

结构记忆还意味着，撤回不等于恢复。一个新基础设施会改变土地使用，一个新平台会改变用户习惯，一个新标准会改变供应链，一个新算法会改变人的预期，一个新制度会改变利益分配。即使后来撤回技术，系统也不会完全回到原状：人已经学习了新方式，组织已经围绕新工具调整，资本已经投入，竞争者已经跟进，用户已经形成新期待。所以，技术治理不能只问“出了问题能不能撤”。在高耦合、高后果系统中，撤回并不等于恢复。一个错误的教育评价系统可能已经影响学生路径；一个错误的金融模型可能已经改变资产配置；一个错误的推荐机制可能已经改变公共讨论结构；一个仓促部署的医疗AI可能已经改变医生和患者之间的信任。系统改变以后，会留下记忆。

这种不可逆性不只是物理不可逆，也包括制度不可逆、行为不可逆、认知不可逆和信任不可逆。桥梁倒塌不可逆，药物副作用不可逆，这些很容易理解；但信任受损、标准锁定、路径依赖、组织惯性，也具有不可逆或半不可逆性质。它们未必违反物理意义上的可回退，却会让社会意义上的恢复极其困难。高后果系统需要更谨慎的技术接入。不是因为技术不能试错，而是因为有些试错会改变系统记忆：低风险比特任务可以快速重试，因为错误容易清除；高风险制度系统不能无限高频试错，因为错误会沉积为新的结构。医疗、教育、金融、法律、交通、能源和公共治理，都需要在试验与不可逆后果之间保持边界。

基础设施的电动化给出了这类不可逆的规模。把一支车队从燃油换成电动，表面上是换车，实际上要同时改造停车场供电、配电容量、消防规范、维修体系、电池回收链条和保险条款；一旦改造完成，逆向切换并不只是把车换回来，因为配电容量、场地布局和维修能力已经按新用途定型。这类改造的时间尺度由施工能力和社会协调速度决定，而不由车辆技术决定，因此它的推进是线性的、可预测的，也很难被一次决策整体加速。

\subhead{整理历史}

结构记忆也提示我们，真正的技术进步不是简单替换旧系统，而是处理旧系统。一个新技术如果不能兼容历史路径，就会遭遇强阻力；如果只是完全服从旧路径，又会被旧结构驯化，无法产生真正改变。成熟的技术变革往往发生在两者之间：既理解旧结构为什么存在，又设计新的迁移路径；既不幻想一键重构，也不放弃重构可能。这需要渐进接口、过渡结构、并行系统、数据迁移、人员再训练、责任重新划分和旧制度的退场机制。技术变革最难的地方，往往不是发明新东西，而是让旧东西有秩序地退出。没有退出机制，新技术只能叠加在旧系统上，系统会越来越厚；有了退出机制，技术才可能真正改变结构，而不是制造更多层级。

结构记忆并不只是负担。没有结构记忆，系统就没有稳定性。医院流程、法律制度、工程标准、交通规则、金融信用、科学共同体规范，都依靠记忆维持可预期性，过去经验被固化为制度，系统才能减少重复试错。问题在于，记忆既保存经验，也保存错误；既提供稳定，也限制变化。它让系统能够运行，也让系统不容易重新开始。成熟系统不是没有历史，而是能够整理历史：它知道哪些旧接口应该保留，哪些旧标准必须退出；哪些经验是安全边界，哪些经验只是过时习惯；哪些沉没结构必须维护，哪些沉没结构应该逐步松动。一个完全失忆的系统没有稳定性，一个无法整理记忆的系统没有可塑性。

AI时代尤其需要这种意识。AI很擅长生成未来方案，却不能自动消除历史负债：它可以告诉我们一个系统应该如何运行，却不能替我们清理旧数据、旧标准、旧利益、旧习惯、旧责任和旧信任结构。未来不是从空白处开始的，每一个新系统，都必须进入一个已经被过去写满的现实。复杂系统有记忆，它不是沿着抽象最优路径运动，而是沿着已经被历史投资、制度接口和行为习惯塑形的路径运动。AI可以生成未来状态，但未来必须穿过过去留下的结构。结构记忆不是过去的记录，它是过去留在系统中的残余、迟滞和不可逆。

结构记忆解释了系统为什么不能从零点出发，也解释了为什么回退本身是一项工程。它还留下一个更基本的问题：在一个已经被历史写过的系统里，究竟哪些状态是真正可以抵达并维持的？记忆让一部分方向变便宜，也让一部分方向变昂贵，可达范围由此被压缩。哪些状态还能到达，答案不在表示空间里。

\subsection{可达空间：局部最优与布雷斯悖论}

AI可以扩展可想象空间，但现实系统真正需要的是可达空间。前面几节讨论了连接、耦合密度、阈值、历史和层级过滤如何限制系统改变，它们最终指向同一个结论：现实系统的状态空间不是完全开放的，一个目标可以被想象、计算、模拟和优化，但它未必能够被系统稳定抵达。这里必须区分两个概念：可想象状态和可达空间。

\subhead{可想象状态与可达状态}

可想象状态，是表示空间中的状态，它可以是一张规划图，一个最优解，一个模型输出，一个制度设计，一个工程方案，一个未来蓝图；AI极大扩展了这种空间，它可以同时生成许多路线、许多结构、许多分子、许多流程、许多组织方案，它让“还有没有别的可能”变得更容易回答。可达状态，则是现实系统在给定约束下真正能够抵达并维持的状态，目标本身成立只是必要条件，此外还要从当前状态出发，存在一条资源可承受、时间可展开、结构可吸收、风险可控制、失败可处理的路径。可想象状态关心终点能不能被描述；可达状态关心从这里到那里能不能走过去。

\key{技术乐观主义最常见的误判，是把可想象状态误认为可达状态。}它看到AI能生成更优目标，就以为现实会自然向那个目标移动。但现实系统不是在空白状态空间中自由漂移，它像一个带有坡度、沟壑、旧路、桥梁、断崖和沉没结构的地形：目标可以在地图上标出，真正困难的是有没有路，路是否太陡，桥是否承重，途中是否有补给，失败后能否撤回。可想象状态的数量由生成能力决定，可达状态的数量由介质、结构和历史共同决定，表示空间中的完备并不等于现实中的可达。

技术落地的关键不是目标本身，而是中间态。很多方案之所以悬浮，不是因为终点错误，而是因为它没有设计中间态：它告诉系统应该成为B，却没有说明系统如何从A走向B。现实改变很少是一次跳跃，而是一串连续或半连续的过渡，旧结构不能瞬间消失，新结构不能瞬间承担全部功能，系统需要并行期、试验区、兼容层、缓冲带、替代路径和退出机制；没有这些中间态，目标越完美，越像悬在空中的图像。可达空间会被多种因素压缩：资源限制决定哪些状态承担不起，时间限制决定哪些路径来不及，耦合限制决定哪些局部变化无法单独发生，历史限制决定系统不能从零点出发，风险边界则排除那些失败后果不可接受的路径。现实系统真正能够选择的，不是所有可能状态，而是被这些条件共同筛选后的少数路径。

中间态的设计有它自己的检验标准：它必须让系统在过渡期内保持可以运行，也必须让过渡期本身有退出的方式。并行运行两套系统要付出双份维护成本，过渡期越长，成本越高；过渡期太短，旧流程尚未完成交接，新流程就会出现功能缺口。可行的设计通常把过渡期按功能切段，让每一段都能独立验收和独立回退，而不是把整个系统当成一个开关。把可达性当成路径问题来处理，就会得到这类具体的设计问题；把它当成目标问题来处理，只会得到越来越多的方案。

\subhead{布雷斯悖论与多主体均衡}

可行路径与实际均衡需要分开。允许忽略新道路时，新增道路不会缩小原来的可行路径集合；但它可能改变司机的选择，令用户均衡中的总通行时间上升。布雷斯悖论处理的正是这种偏离。新增工具、指标和模型是否产生类似效果，要说明相应行动者与成本结构。

布雷斯悖论的意义不在交通，而在可达性的定义。它说明，路径成本不是路径本身的静态属性，而是所有主体共同使用之后生成的结果；路径不是单纯的几何通道，而是拓扑结构、成本函数、主体选择和流量分布共同形成的动态关系。新增连接改变的不是地图上的一条线，而是系统中的成本地形；成本地形一变，行为分布就会重排；行为分布一重排，那条看似更短的路就可能变成新的拥堵源。因此，可达性不是单主体属性，而是多主体均衡属性。

这个性质可以用用户均衡条件写出来。设 $\mathcal{P}_w$ 为起讫对 $w$ 之间所有可行路径的集合，$x_f$ 为路径 $f$ 上的流量，$\mathbf{x}$ 为全体路径流量构成的分布，$c_f(\mathbf{x})$ 为路径 $f$ 的通行成本，其数值由全体流量共同决定，单位是时间。用户均衡是这样一种分布 $\mathbf{x}^{*}$：对每一对起讫 $w$，凡是流量为正的路径，其成本都不高于同一对之间任何其他路径的成本，
\begin{equation}
c_f(\mathbf{x}^{*}) \le c_g(\mathbf{x}^{*}),
\qquad \forall f\in\mathcal{P}_w \text{ 且 } x_f^{*}>0,\ \forall g\in\mathcal{P}_w .
\end{equation}
这个条件刻画的是个体最优与系统最优的偏离：每个司机只比较自己的成本，没有人对系统总成本负责，因此均衡可以停在一个所有人都想离开、但没有人能单独改善的状态。条件成立的前提是出行者信息充分、路径可以自由选择，并且只按自身通行时间决策；一旦收费、限行或信息干预介入，均衡位置就会移动。一个人可以走通，不代表系统可以走通；一条路对个人是低阻，不代表对系统是低阻；一个方案对局部是更优，不代表在共同使用之后仍然更优。复杂系统的可达空间，不能只看“有哪些路径”，还要看这些路径被所有主体同时使用后，系统会落到什么状态\cite{wardrop1952some}。

个体最优与系统最优之间的差距可以被量化，结果比直觉更严重。在延迟随流量线性上升的网络里，自私路由所形成的均衡，其总通行时间最多是系统最优解的三分之二到四分之三区间内的某个倍数，也就是说，即使所有出行者都做出了对自己最优的选择，系统整体的时间损失也有一个不到 1.4 倍的确定上界，而某些网络上这个上界会被真正达到\cite{roughgarden2002how}。更值得注意的经验证据来自道路供给本身：对美国城市数据的研究发现，一个地区的高速公路车道里程增加约 10\%，长期看车辆行驶里程也会增加约 10\%，通行时间的改善远小于按静态流量估算的结果\cite{duranton2011fundamental}。新增容量改变了出行者的选择，包括路径、出行时段和居住地点的选择，车流因此不是被消化，而是被重新生成。

\subhead{路径膨胀与均衡检验}

这对AI时代尤其重要。AI会不断增加“路径”：更多推荐路线，更多生成方案，更多组织工具，更多评价指标，更多自动化决策，更多低成本入口，它会让许多原本昂贵、低频、谨慎使用的选择变得便宜、高频、普遍。但一条低阻路径一旦被大规模使用，就不再只是低阻路径，它会吸引流量，诱导模仿，制造拥堵，产生套利，推动同质化，最后改变系统均衡。导航可以让所有人涌向同一条小路，推荐算法可以让创作者集体迎合相同信号，科研工具可以让大量研究同时挤向同一批候选方向，组织指标可以让不同部门围绕同一套数字变形；局部看，每个主体都在走更优路径，整体看，系统可能进入更窄、更拥挤、更脆弱的状态。可达空间不是路径数量，而是路径被共同使用后仍能维持的状态集合。

这也要求重新理解冗余、备用路径和缓冲。冗余不是随意增加连接，备用路径也不是路径数量越多越好。布雷斯悖论提醒我们，某些新增连接会诱导坏均衡；真正有价值的冗余，是能增加系统的可选择性、可隔离性、可恢复性和可回滚性的结构。它不是让所有主体同时涌向同一条新捷径，而是在扰动出现时提供分流、隔离、等待、撤回和重组的余地。换言之，扩大可达空间的不是任意路径，而是\term{可治理路径}。可治理路径有几个特征：它不会把所有流量吸向同一瓶颈；它能在异常时分担压力；它能隔离局部失败；它能保留回滚余地；它不会把系统推入更差的均衡。相反，那些看似低阻、实际会诱导所有主体同步涌入的捷径，可能不是冗余，而是新的拥堵入口。

电力系统把冗余写成了一条明确的准则。规划与运行中普遍采用的做法是，在任一主要元件故障退出之后，系统仍要保持稳定运行并不损失负荷，重要断面还要在此基础上再承受一次故障，这就是通常所说的 N-1 与 N-2 检验\cite{morison2004power}。这条准则的代价是实打实的：需要预留备用容量、保留输电走廊、维持一定数量的机组处于可调状态，日常运行的经济性因此低于理论最优。它之所以被长期保留，是因为电网的失败后果不可回滚，备用容量买的不是效率，而是系统在扰动之后仍然能回到安全状态的能力。这正是可治理冗余与拥堵入口的分界：前者在异常时才被使用，在常态下不改变主体的选择；后者在常态下就吸引所有流量，把系统推向更窄的均衡。

\subhead{局部最优与可治理路径}

可治理路径之所以难以打开，是因为系统往往停在局部最优。所谓局部最优，不是一个系统真的达到了最好状态，而是它停在了附近最容易维持的状态。这个状态可能低效，可能笨重，可能不公平，可能已经过时，但它稳定。稳定本身就是一种力量：一个流程一旦被组织适应，一个标准一旦被产业采用，一个评价体系一旦被行为围绕，一个城市格局一旦被生活嵌入，它就不再只是一个选择，而变成了环境。局部最优的本质，不是“最好”，而是“离开它很难”。很多系统不是不知道更优方案，而是无法承受离开当前状态的代价。一个旧软件系统不好用，但所有数据、权限、流程和培训都围绕它展开；一个评价机制有缺陷，但学校、教师、学生和资源分配已经围绕它运行；一种行业标准并不完美，但供应商、设备、认证和客户都依赖它。它们不是全局最优，却是可维持的局部盆地。

AI会让这个问题更加明显。它可以揭示更多可能的高处，也可以让人更清楚地看到当前系统的不完美。但看见高处，不等于拥有上山的路；没有路径设计的智能，只会制造更清楚的焦虑：知道哪里更好，却不知道怎么过去。所以，AI在复杂系统中的高阶价值，不只是生成目标，而是帮助发现可达路径，它应该帮助系统识别哪些路径可以分阶段推进，哪些接口需要先改，哪些旧结构必须暂时保留，哪些节点可以先试验，哪些风险必须隔离，哪些失败必须可回滚。但路径设计不能只问单个路径是否更短，还要问路径开放之后会改变什么均衡。一个新通道如果只是降低局部成本，却诱导所有主体涌向同一处，它可能缩小系统可达空间；一个新工具如果只是让所有人追逐同一个指标，它可能强化局部最优；一个新模型如果只是让所有决策趋向同质化，它可能提高短期效率，却降低系统适应性。路径设计真正要处理的，不是“有没有路”，而是“这条路被使用之后，会不会把系统推向更窄的状态”。

\subhead{效率与鲁棒性}

可达空间还涉及鲁棒性。除了均衡，它还取决于系统在扰动出现时能否继续维持这些路径可用。很多系统为了追求效率，会不断减少冗余、压缩缓冲、提高利用率、细化控制、增加连接、消除备用路径。短期看，这会让系统更快、更便宜、更精确；长期看，它会缩小可达空间。一个没有冗余的系统，在正常条件下可能表现优异，但遇到扰动时缺少替代路径，它能沿着既定轨道高速运行，却很难切换轨道\cite{carlson1999highly}。效率关注在已知环境下把资源用到极致，鲁棒性关注环境变化时系统能否不崩溃；效率常常要求减少闲置，鲁棒性常常要求保留余量；效率喜欢紧密连接，鲁棒性需要隔离和缓冲。效率最大化的系统，在正常时期可能非常漂亮，但一旦外部条件改变，就会暴露出路径贫乏的问题。

医院排班如果被优化到极限，平时利用率很高，但突发事件来临时缺少缓冲；电网调度如果只追求短期成本最小，面对极端天气和负荷波动时可能缺少弹性；平台算法如果只追求点击率，可能牺牲长期信任和公共讨论质量；组织如果只追求即时产出，可能没有时间学习、修复和转型。冗余不是简单浪费，缓冲不是简单低效，备用路径不是落后。它们都是可达空间的一部分：没有冗余，系统只有一条路；没有缓冲，系统无法等待；没有备用路径，系统无法绕行；没有回滚机制，系统不敢试验。一个完全被压缩到极致效率的系统，往往也是最难改变、最难恢复、最难吸收新技术的系统。但冗余必须是可治理冗余，它不是盲目增加连接，也不是无限增加选项。真正的冗余，是在异常时打开，在常态下不破坏系统均衡的余地；它不是让系统永远更复杂，而是让系统在关键时刻不被单一路径锁死。这里的判断标准仍然是均衡：一项冗余如果只是给所有主体提供了一条共同的新捷径，它就不是冗余，而是新的拥堵入口。

\subhead{可塑性}

因此，可塑性比效率更接近长期生命力。系统可塑性，是一个系统在不崩溃的前提下改变自身结构的能力，它不是松散，也不是混乱，而是保留调整余地。一个系统如果太松散，就不能形成有效功能；如果太紧密，就不能适应新环境。真正成熟的结构，应该既能运行，又能改造；既有秩序，又有余地；既能利用局部优化，又不被局部优化锁死。

从这个角度看，技术进步不只是让系统更快抵达某个目标，而是扩大系统的可达空间：降低资源成本，可以扩大可达空间；缩短必要过程，能让更多路径来得及被走到；减少不必要耦合，把可达域向外推；处理历史负债，让原本走不通的路径变得可行；建立沙盒、缓冲和回滚机制，也可以扩大可达空间。但所有这些都必须接受均衡检验：它们是否真的增加可承受路径，还是只是制造新的低阻拥堵点。这也是前面几节的讨论最终指向的含义：局部优化不等于系统改变，因为系统需要路径；耦合密度会让节点变重，因为连接会降低自由度；凝胶化会让系统突然变硬，因为连接越过阈值后局部扰动会贯通全局；结构记忆会让系统不可任意重置，因为过去会塑造未来；布雷斯悖论进一步说明，新增路径本身也可能改变均衡，使可达空间反而收缩。

AI可以在表示空间中打开极大的可能性。但如果现实系统的可达空间没有扩大，大量可能性就会停留在模型、报告、试点、规划和愿景中，它们存在于可想象空间，却不能成为系统状态。所以，未来真正稀缺的不是目标，而是路径；不是想象力，而是可达性；不是全局最优的图像，而是从局部最优走出去的通道。管理黏滞，最终就是管理可达空间：它不是否认目标，也不是放弃优化，而是把目标放回现实路径中考察，哪些阻力可以降低，哪些连接可以松动，哪些新增路径会诱发坏均衡，哪些冗余必须保留，哪些中间态必须设计，哪些试验必须隔离，哪些失败必须可回滚，哪些旧结构必须有秩序地退出。AI扩展的是可能性的边界；现实需要的是路径的打开。但路径不是越多越好，真正的可达空间，不是地图上画出了多少路，而是系统在所有主体共同选择之后，仍然能够稳定抵达并维持多少状态。

到这里，可达空间一直被当作结构问题来处理：连接方式、历史残余和均衡分布决定哪些状态可以被抵达。但还差一层没有计入。一条路径能不能走通，不只取决于它在地形上的位置，也不只取决于有多少主体同时使用它，还取决于路径两端各挂着多少关系。同一个改动，落在一个孤立节点上只是一次改动，落在一个被依赖网络包围的节点上，就变成一次结构移动。

\subsection{结构的重量：依赖面、同步成本与等效惯性}

路径能不能走通，只回答了系统可以到达哪里，没有回答移动的代价从哪里来。同样是替换一个部件，有的系统里它是一次维护，有的系统里它是一次重新设计；同样是上线一个模型，有的场景里它是一次试验，有的场景里它是一次责任重组。这个差别不能由路径长度、连接数量或均衡位置单独解释。把这种依赖拆开，需要三个可以分别丈量的变量：一个改变会牵动多少关系，这些关系是否必须同时移动，失败之后能不能退回。三者同时上升时，系统表现出来的就是\term{等效惯性}。

\subhead{等效惯性}

但移动的代价不只取决于路径的长短。一个孤立部件很容易移动；一个嵌入网络的部件，即使自身很小，也可能很难移动，因为改变它，不只是改变它本身，还要改变它与周围节点之间的关系。一个独立脚本可以很快修改，一个接入财务、客户、库存、审批和合规流程的软件模块就很难修改；一个实验室模型可以快速更新，一个进入医院、银行、电网或交通系统的模型就不能随意更新，它不再只是模型，而是系统关系中的一个节点。

这就是结构的重量。前面讲网络耦合密度时，把附加在节点上的改动成本记作虚拟质量，那里只算了一项：改动会牵动多少关系；把同步要求和回滚代价一并计入，同一个量才展开成完整的三项。\key{结构的重量不是物理质量，而是关系赋予系统的等效惯性。}一个节点连接越多，它的变化越不再属于自己：它必须同时满足更多接口、更多规则、更多上下游依赖和更多责任边界；它看起来只是一个点，实际拖着一片结构。更准确地说，结构重量来自协同移动成本。改变一个节点时，有多少关系必须一起调整，这些关系是否必须同步调整，调整失败后能否局部回滚，这三件事决定了一个结构的等效惯性。部件本身可能很轻，但它的依赖面越大，同步要求越高，回滚越困难，它就越重。

三个变量很少各自单独起作用。它们同时上升的时候，代价不是三份，而是一次整体移动：依赖面决定了要动谁，同步成本决定了能不能分开动，回滚代价决定了动了之后还退不退得回来。写一个内部数据处理脚本，改一行、当天上线、出错就恢复，三个变量都接近于零；换掉医院里一条计价规则，要同时改动收费系统、医嘱系统、医保接口、对账流程和患者告知，任何一个环节没跟上，账目就对不上。同一个技术能力，装在前一个位置上是工具，装在后一个位置上是一项工程。这也解释了为什么同一个团队在两种位置上会表现出完全不同的效率：能力没有变化，被拖动的结构变了。

\subhead{依赖面}

三个变量当中，最容易度量的是依赖面。\term{依赖面}，是一个改变会牵动多少关系。在低依赖面场景中，局部改变可以保持局部性：改一段个人笔记，换一张广告图，修改一个独立脚本，通常不会牵动复杂关系；失败可以回滚，成本可以隔离，责任也相对简单。这样的场景对AI很友好，因为模型输出可以快速试错，局部效果几乎可以直接转化为实际效果。在高依赖面场景中，局部改变不再局部：医院里的一个诊断模型，会牵动医生流程、病人信任、医保支付、责任认定和监管规则；银行里的一个风控模型，会影响客户授信、资金配置、合规审查和风险暴露；电网里的一个调度系统，会连接发电、负荷、储能、安全和事故后果。它们不是单点工具，而是依赖网络中的节点。

反过来看，技术对象自身可能很轻，但它所拖动的结构很重。一个模型文件可以快速替换，但如果它已经嵌入业务规则、人员分工、客户承诺、合规流程和责任体系，它就不能像文件一样自由替换；一个界面按钮可以瞬间改变，但如果它改变用户行为、交易路径和风险分布，它就具有结构后果。依赖面不是抽象的规模，它可以直接换算成可替换性：能单独替换的对象才算轻，必须连带替换的对象就变重了。

依赖面在工程上有可以数出来的对应物：一次改动要触及多少份接口文档，经过多少个审批环节，涉及多少个责任主体，牵动多少家外部供应商。工业系统里有一个物理版本。一座大型高炉从点火到停炉大修，炉役通常在十五到二十年之间，期间不能停机，任何涉及炉体、送风、冷却和原料系统的改造，都必须安排在短暂的休风窗口内完成\cite{dong2022operational}。这不是技术难度的问题，而是依赖面的问题：一座高炉连着一整条钢铁流程的物流、能源和订单，它不是一个可以随时替换的模块。同样地，一条城市主干道、一座变电站、一段铁轨，它们自身的建设周期并不离谱，真正长的是它们所嵌入的既有网络完成重新适配所需的时间。

\subhead{同步成本与回滚代价}

依赖面之外，还有第二个变量：这些关系是否必须同时移动。\term{同步成本}，是这些依赖关系是否必须同时移动。有些系统可以分段更新，局部试验，逐步迁移，它们的同步成本低：一个模块可以先替换，一个接口可以先兼容，一个新流程可以先在边缘场景运行；即使失败，影响也被限制在局部。这样的系统更容易吸收新技术。另一些系统必须整体协同：标准要一致，数据要同步，责任要重新划分，人员要一起培训，旧系统和新系统要并行兼容。一个环节改了，其他环节不改，系统就出现断裂。同步成本越高，结构越重。很多技术不是不能用，而是不能被系统同时挪动。

第三个变量是失败之后的退路。\term{回滚代价}，是调整失败后能否退回。低回滚代价的系统，可以大胆试：软件功能、内容推荐、局部流程、低风险辅助工具，都可以灰度、沙盒、回滚。高回滚代价的系统，则不能轻易试。医疗诊断、金融风控、交通调度、电力控制、公共治理，一旦改变流程和责任，失败就不是简单恢复旧版本。人已经学习新方式，组织已经调整资源，用户已经形成预期，风险已经发生转移。依赖面、同步成本和回滚代价叠加起来，就形成等效惯性。

同步成本有一个古老而干净的版本，出现在度量单位上。1999 年，火星气候探测者在进入火星大气时失联，事后调查认定，地面导航软件的一组推力数据以磅力秒为单位输出，而探测器上的推进系统按牛顿秒解释，两条链路各自内部一致，接口处却缺少共同单位，误差在长时间的轨道修正中累积成一次不可逆的进入角偏差\cite{nasa1999mishap}。失效的原因不在某一层算错了，而在两层没有在同一时刻使用同一套约定。单位、字段、时间戳、坐标约定、状态定义，这些看起来琐碎的约定，恰恰是依赖面最小的接口；它们的数量随系统复杂度增长，任何一个没对齐，同步成本就会在联调阶段集中暴露，而联调阶段通常已经无法回退到设计阶段。

回滚代价最容易被低估的地方在于，有些中间状态根本无法保存。银行核心账务系统的迁移是典型场景。2018 年英国 TSB 银行把客户数据迁移到新平台，迁移过程中出现大规模账户访问故障，约一百九十万客户受影响，故障持续数周，事后成本超过三亿英镑\cite{tsb2019migration}。麻烦在于，数据一旦写入新系统并产生新的交易，旧系统就不再持有完整状态，回退只能靠第二次迁移完成，而第二次迁移面对的是一个已经改变的数据集。这也是高回滚代价系统的常见做法：把新方案先放在旁路位置，先并行、再切换，切换点之前的每一步都可撤，切换点之后的每一步都不可撤。

耦合还有第三种后果：单条输入可能直接获得系统权限。波音 737 MAX 的机动特性增强系统只依据单个迎角传感器的读数触发，在两起事故中，传感器读数异常时系统反复压低机头，2018 年和 2019 年的两起坠机造成 346 人遇难，全球机队停飞接近二十个月\cite{jatr2019boeing}。事后改动的方向说明了问题所在：改的不是算法的聪明程度，而是让判断依赖两个传感器的比较，并在不一致时把控制权交还飞行员。冗余在这里划出的是一条权限边界：哪些输入可以直接获得执行权限，哪些必须经过比较和确认。依赖面、同步成本、回滚代价三者都可以被设计改变，但改变它们的手段通常不是更强的算法，而是结构上的隔离、并行与可撤销。

\subhead{扩散差异与隐藏的重量}

一个反复出现的格局是：某些看似强大的AI应用长期停留在试点，而某些看似普通的AI应用反而迅速扩散。文档处理、客服初筛、代码补全、合同条款抽取、局部质检、图像标注，往往更容易落地。它们不是最宏大的应用，但边界清楚，输入输出明确，失败容易隔离，责任相对可控；它们嵌入位置轻，所以扩散快。相反，进入核心流程的AI更难。它可能更有价值，却拖着更大的依赖面：人员分工、系统接口、合规要求、责任归属、用户信任、长期维护。技术能力只是其中一部分，真正困难的是让整个依赖网络一起移动。

这条规律在企业统计里也能看到形状。AI应用往往先在文本和图像的边缘位置聚集：合同条款抽取、工单分类、代码补全、巡检图像标注，这些位置的共同点是输入输出可以写成一张明确的清单，错误可以被下游的人工复核接住。一旦进入授信、诊断、调度、定价这类位置，同样的模型就要接入责任链条，需要回答它依据什么、由谁复核、出错由谁承担，落地周期随之从月变成年。位置决定了落地速度，而位置可以在部署之前被选择，这是少数可以由技术团队自主决定的变量。

更麻烦的是，结构重量常常隐藏在顺滑体验背后。用户端越顺滑，后台往往越复杂。一次点击背后，可能是支付系统、身份认证、风控模型、物流网络、客服体系、数据合规和平台规则的协同。用户看到的是低摩擦界面，系统承担的是高耦合维护。技术把摩擦从显眼处移走，却不一定让摩擦消失；它可能只是把摩擦转移到后台结构中\cite{star1999layers}。衡量一个系统是否容易被改造，不能看它暴露给用户的界面，而要看它背后有多少条依赖必须同时成立。

支付清算给出了一个具体的形状。每天交易结束后，系统必须在夜间的一个固定窗口内完成对账、差错处理与资金归集，白天的业务才能在新的余额上开始；这个窗口对用户完全不可见，却限制了白天哪些参数不能随意改动，也限制了任何一个上游系统的发布时间\cite{kahn2009why}。用户端的顺滑与后台窗口的刚性同时存在，只是后者不进入任何一份产品说明。评估一项改动的真实代价，必须把这类不进入界面、但决定发布时间表的环节算进来。

\subhead{连接的双重代价与模块化}

因此，连接不是单纯的好事。连接让系统能够协同，也让系统获得惯性。没有连接，系统只是松散部件，无法形成复杂功能；连接太密，系统又会变得沉重，局部调整很难不牵动全局。现代社会的效率，很大程度上来自高度耦合：电网、金融网络、医疗体系、供应链、互联网平台、交通系统和科研体系，都依靠密集连接获得规模、速度和协同。但同样因为密集连接，它们也失去了许多局部自由。

这是一种基本交换：系统用自由度换取协同能力。好的结构不是连接最多的结构，而是连接被恰当地组织。关键部分需要稳定连接，非关键部分需要保持可替换；模块内部可以紧密耦合，模块之间要有清晰接口；系统可以提高效率，但必须保留缓冲、冗余和回滚空间。结构不是越紧越好，也不是越松越好，真正重要的是控制耦合的位置和强度\cite{boccaletti2006complex}。

模块化的意义就在这里。模块化不是把系统机械拆碎，而是控制依赖面。它让局部变化不必无限扩散，让失败可以被隔离，让新技术有可插入的位置。标准化接口、清晰输入输出、责任边界、沙盒环境、灰度部署、可回滚机制，都是降低等效惯性的方式。它们不是让系统无结构，而是让结构可管理\cite{baldwin2000design}。这也区分了结构重量与拓扑耦合：拓扑耦合说明局部为什么会牵动全局，结构重量说明被牵动的全局为什么难以移动。

模块化在产业史上有过一次彻底的示范。集装箱出现之前，货物靠人力逐件搬运，一艘货轮的装卸以周计，港口费用在运输成本中占很高比例；标准化的钢制箱体把装卸变成吊装，货物在工厂装进箱子，中途不再拆开，运费结构随之改变\cite{bernhofen2016estimating}。这件事的技术含量不高，它的作用是切断了货物与船舶之间的依赖面：箱子承担适配，船和货不再互相牵制。标准化的作用是把变化集中到一个可以单独替换的层上，箱体承担适配，船舶与货物各自演化。

模块化也有失败的方向。把单体系统拆成微服务，接口数量随服务数量增长，一次业务改动需要多个服务同时发布；如果拆分只发生在代码层面，而没有把数据、事务和发布节奏一起切开，耦合就从系统内部搬到了网络接口上，运维成本反而上升，工程界给这类结构起了一个名字，叫分布式单体\cite{taibi2018architectural}。模块化成不成功，判据是依赖面有没有真的变小：如果一次改动仍然要求多个部分同时上线，拆分只是换了耦合的位置，没有改变结构重量。

\subhead{智能与结构的匹配}

AI要真正进入现实系统，必须面对这种结构重量。一个AI工具如果没有清晰输入、清晰输出、适用边界、失败处理和责任归属，就很难规模化部署。不是因为模型不够聪明，而是因为它进入了高依赖面、高同步成本、高回滚代价的位置。相反，一个边界清楚、风险可隔离、失败可回滚的场景，即使不够宏大，也更容易成为真实能力。技术落地不是单纯比较智能水平，而是比较智能与结构的匹配程度。

匹配可以写得更具体。同样的模型能力，放在辅助位置上，输出由人复核，错误可以被拦截，依赖面小；放在决策位置上，输出直接进入执行，错误沿依赖网络传播，依赖面大。前一种位置的成功经验不能直接搬到后一种位置，结构位置不同，同一套输出进入的后果链条也不同。选择部署位置时要回答的是三个问题：这个输出会触发哪些动作，这些动作能否被隔离，被隔离之后系统能否继续运行。

反过来看，AI可以让一个节点更聪明，但如果这个节点拖着整张网络，它仍然移动得很慢。结构的重量，不在节点本身，而在节点必须一起带动的关系。改变一个点很容易，改变一个被关系包围的点很难。结构越重，技术越不能只问“能不能做”，还要问：依赖面多大，同步成本多高，失败能不能退回。

依赖面、同步成本和回滚代价可以逐个丈量，但它们的叠加并不总是平滑累积。这些成本同时上升之后，系统先失去的是局部性：一个环节的调整会把整片结构一起拉进来，一个失败不再容易被隔离在发生地。到这一步，问题已经从“有多重”变成“是什么状态”，而后者不是靠逐个丈量能够回答的。三个变量的叠加还会改变介质的响应方式：同样的扰动，落在低耦合结构上会被吸收，落在高耦合结构上会被放大。要描述这种差别，需要的不再是一组独立的成本项，而是一个能同时容纳资源、时间和结构的位置坐标。
